% !TEX program = xelatex
% ============================================================================
% Planteamiento del problema, supuestos, reglas de negocio y diseno
% Proyecto Final de Bases de Datos - Inventario, Compras y Ventas (PostgreSQL)
% Entrega 2 de 4 - Sincronizado con Scripts sql/
% ============================================================================
\documentclass[11pt,oneside]{article}

\usepackage{estilo-proyecto}
\docnombre{Planteamiento y dise\~no del sistema}

\addto\captionsspanish{\renewcommand{\contentsname}{Contenido}}
\setcounter{tocdepth}{2}
\setlength{\emergencystretch}{1.6em}

% Notacion compacta de esquema relacional (raggedright para evitar desbordes)
\newcommand{\rel}[2]{%
  \par\vspace{0.28em}%
  {\raggedright\noindent
   {\ttfamily\small\bfseries\color{acento}#1}\enspace$=$%
   \enspace({\ttfamily\small #2})\par}}
\newcommand{\rpk}[1]{\underline{\textbf{#1}}}
\newcommand{\rfk}[1]{\textcolor{fk}{#1}\textsuperscript{\scriptsize\sffamily FK}}
\newcommand{\ruq}[1]{#1\textsuperscript{\scriptsize\sffamily UQ}}

\begin{document}

\portadilla{Proyecto final \textperiodcentered{} Bases de datos \textperiodcentered{} Entrega 2 de 4}
  {Planteamiento del problema y dise\~no arquitect\'onico}
  {Alcance, supuestos, reglas de negocio, fragmentaci\'on y modelo relacional (3FN)}

\tableofcontents
\vspace{0.5em}

% ============================================================================
\section{Planteamiento del problema}

Una empresa comercial opera m\'ultiples sedes de venta al p\'ublico apoyadas por
un centro de distribuci\'on principal. Cada sede gestiona almacenes o bodegas
propias, administra inventarios por producto, registra pedidos de venta,
recauda pagos, tramita devoluciones y coordina tanto compras a proveedores como
transferencias internas de mercanc\'ia entre bodegas.

A medida que crecen las sedes y el volumen transaccional, centralizar todas las
operaciones en un esquema monol\'itico sin separaci\'on de responsabilidades
dificulta el control de accesos por \'area, mezcla datos maestros con
operaciones diarias e impide aislar la informaci\'on por sede. Al mismo tiempo,
la direcci\'on general y el equipo de auditor\'ia requieren reportes
corporativos consolidados (valorizaci\'on de existencias, alertas de
reabastecimiento, saldos por cobrar, cumplimiento de proveedores y balance
por sede) sin duplicar conteos ni perder la trazabilidad hist\'orica de
precios y costos.

Para resolver esta necesidad, se dise\~na e implementa en PostgreSQL una capa
de datos estructurada bajo una \textbf{arquitectura l\'ogica de microservicios
por esquema}, combinada con una estrategia de \textbf{fragmentaci\'on horizontal
por sede} sobre los conjuntos operacionales de inventarios y ventas, y una
capa dedicada de \textbf{vistas de solo lectura} para integraci\'on corporativa.

% ============================================================================
\section{Alcance y condiciones m\'inimas}

\subsection{Alcance funcional del proyecto}

El sistema cubre el ciclo completo de datos maestros, abastecimiento,
existencias, ventas y control de acceso en PostgreSQL, organizado en cuatro
entregas progresivas. Hasta la \textbf{Entrega 2 de 4}, el alcance cubierto
y proyectado comprende:

\begin{itemize}[leftmargin=1.5em, itemsep=3pt]
  \item \textbf{Multi-sede y multi-almac\'en:} soporte m\'inimo para cuatro
        sedes (incluida una de distribuci\'on) y cuatro bodegas; la propuesta
        delimitada implementa \textbf{5 sedes} (1 de distribuci\'on, 3 de
        ventas y 1 regional) y \textbf{8 almacenes}.
  \item \textbf{Modularidad por microservicios:} separaci\'on expl\'icita en
        cinco esquemas transaccionales (\texttt{catalogo\_\allowbreak terceros},
        \texttt{compras\_\allowbreak abastecimiento},
        \texttt{inventarios\_\allowbreak transferencias},
        \texttt{ventas\_\allowbreak pagos},
        \texttt{identidad\_\allowbreak accesos}) y un esquema
        exclusivo de lectura (\texttt{reportes}).
  \item \textbf{Modelado e integridad:} diagrama conceptual Entidad--Relaci\'on
        en notaci\'on de Chen extendida con agregaciones (\texttt{A1}--\texttt{A3}),
        esquema relacional normalizado hasta 3FN, diccionario de datos y
        restricciones \texttt{PRIMARY KEY}, \texttt{FOREIGN KEY},
        \texttt{UNIQUE}, \texttt{NOT NULL}, \texttt{DEFAULT} y \texttt{CHECK}.
  \item \textbf{Poblaci\'on reproducible:} carga determin\'istica con al menos
        cuatro meses de hist\'orico (semilla \texttt{setseed(0.42)}) que cubre 100
        productos, 6 categor\'ias, 12 proveedores, 120 clientes, 25 empleados
        y usuarios, 75 \'ordenes de compra, 350 pedidos de venta, 30
        transferencias y 25 devoluciones.
  \item \textbf{Distribuci\'on, vistas y anal\'itica (proyecci\'on Entregas 3--4):}
        fragmentaci\'on horizontal por \texttt{id\_sede} sobre inventarios y
        pedidos de venta, 10 vistas SQL consolidadas en \texttt{reportes}, 6
        roles RBAC en PostgreSQL, 20 consultas anal\'iticas avanzadas (CTE,
        funciones de ventana y operaciones de conjuntos) y matriz de pruebas.
\end{itemize}

\subsection{Exclusiones expl\'icitas}

De acuerdo con los lineamientos del curso, el proyecto se concentra en la
ingenier\'ia de datos en PostgreSQL y \textbf{no incluye}:

\begin{itemize}[leftmargin=1.5em, itemsep=2pt]
  \item Programaci\'on de microservicios en lenguajes de aplicaci\'on, API REST,
        mensajer\'ia as\'incrona, despliegue en contenedores Docker ni
        interfaces gr\'aficas web o m\'oviles.
  \item Mecanismos de sincronizaci\'on en tiempo real entre nodos f\'isicos,
        autenticaci\'on web externa ni pol\'iticas de seguridad por fila (RLS).
  \item Pruebas de concurrencia, bloqueos transaccionales, reversi\'on
        distribuida o automatizaci\'on compleja mediante \textit{triggers}.
\end{itemize}

% ============================================================================
\section{Supuestos de dise\~no}

\begin{enumerate}[leftmargin=1.5em, itemsep=4pt]
  \item \textbf{Simulaci\'on acad\'emica en una sola instancia de PostgreSQL.}
        Implementar seis esquemas separados dentro de una \'unica instancia de
        PostgreSQL constituye una \textit{simulaci\'on acad\'emica} de la capa
        de datos de una arquitectura de microservicios. No equivale a desplegar
        microservicios productivos con motores independientes, aislamiento de
        red ni consistencia eventual; su prop\'osito es modelar fronteras claras
        de propiedad sobre las tablas y controlar permisos por esquema.

  \item \textbf{Referencias l\'ogicas frente a claves for\'aneas locales.}
        En el plano conceptual (Diagrama ER), las asociaciones entre tablas de
        distintos microservicios se distinguen como referencias l\'ogicas entre
        esquemas (l\'inea discontinua) para indicar acoplamiento entre servicios.
        No obstante, al ejecutarse todos los esquemas dentro de la misma base
        de datos PostgreSQL, el script
        \texttt{aplicacion\_de\_claves\_\allowbreak foraneas.sql}
        materializa las 32 restricciones \texttt{FOREIGN KEY} (locales e
        inter-esquema) para asegurar integridad referencial estricta durante las
        pruebas.

  \item \textbf{Conservaci\'on hist\'orica de precios y costos.}
        Se asume que el precio de lista de un producto
        (\texttt{productos.\allowbreak precio\_lista}) y el costo de referencia
        del proveedor (\texttt{producto\_\allowbreak proveedor.\allowbreak
        costo\_\allowbreak referencia}) pueden variar en el tiempo. Por ello,
        cada detalle de pedido y de orden de compra captura el
        \texttt{precio\_\allowbreak unitario} y el
        \texttt{costo\_\allowbreak unitario} vigentes al momento de la
        operaci\'on, evitando alterar reportes de periodos anteriores.

  \item \textbf{Moneda \'unica, estados l\'ogicos y alcance fiscal.}
        Todos los valores monetarios se registran en una \'unica moneda
        (\texttt{NUMERIC(12,2)}, sin tabla de tasas de cambio ni discriminaci\'on
        tributaria separada). Las entidades maestras y de personal no se eliminan
        f\'isicamente: emplean borrado l\'ogico mediante \texttt{estado}
        (\texttt{'Activo'} / \texttt{'Inactivo'}) para preservar el historial
        transaccional.

  \item \textbf{Identificadores y granularidad temporal.}
        Las claves primarias son enteros (\texttt{INT}) asignados de manera
        determin\'istica en los scripts de carga reproducible; las fechas
        transaccionales y de bit\'acora se modelan con granularidad diaria
        (\texttt{DATE}). Se asume que un pedido completado admite m\'ultiples
        pagos parciales y eventuales devoluciones asociadas.
\end{enumerate}

% ============================================================================
\newpage
\section{Reglas de negocio}

\subsection{Cat\'alogo y terceros (\texttt{catalogo\_terceros})}
\begin{itemize}[leftmargin=1.5em, itemsep=2.5pt]
  \item \textbf{RN-01 (Sedes y almacenes):} toda sede posee nombre \'unico y
        tipo obligatorio (\texttt{'Ventas'}, \texttt{'Distribuci\'on'} o
        \texttt{'Regional'}). Cada almac\'en pertenece exactamente a una sede y
        se clasifica como \texttt{'Principal'}, \texttt{'Tienda'} o
        \texttt{'Bodega'}.
  \item \textbf{RN-02 (Productos y categor\'ias):} cada producto se identifica
        por un \texttt{sku} \'unico, pertenece a una \'unica categor\'ia (de
        nombre \'unico) y tiene \texttt{precio\_lista >= 0}.
  \item \textbf{RN-03 (Terceros y suministro M:N):} proveedores y clientes
        tienen n\'umero de identificaci\'on o documento \'unico (\texttt{BIGINT}).
        Un producto puede vincularse con uno o varios proveedores en
        \texttt{producto\_proveedor}, registrando \texttt{costo\_referencia >= 0}
        y \texttt{plazo\_entrega >= 0} en d\'ias.
\end{itemize}

\subsection{Compras y abastecimiento (\texttt{compras\_abastecimiento})}
\begin{itemize}[leftmargin=1.5em, itemsep=2.5pt]
  \item \textbf{RN-04 (\'Ordenes de compra):} cada orden referencia un
        proveedor, un almac\'en destino y un empleado responsable; su fecha
        estimada de entrega debe cumplir \texttt{fecha\_estimada >= fecha\_orden}
        y su estado pertenece a \texttt{\{'Pendiente', 'Recibida Parcial',
        'Recibida Total', 'Cancelada'\}}.
  \item \textbf{RN-05 (Detalle y recepciones):} una orden no repite el mismo
        producto en \texttt{detalle\_orden\_compra} (clave compuesta
        \texttt{(id\_orden\_compra, id\_producto)}, agregaci\'on \texttt{A1}) y
        exige \texttt{cantidad\_solicitada > 0} y \texttt{costo\_unitario >= 0}.
        Las recepciones parciales o totales registran \texttt{cantidad\_recibida >= 0}
        por producto en estado \texttt{'Procesado'}, \texttt{'Parcial'} o
        \texttt{'Anulado'}.
\end{itemize}

\subsection{Inventarios y transferencias (\texttt{inventarios\_transferencias})}
\begin{itemize}[leftmargin=1.5em, itemsep=2.5pt]
  \item \textbf{RN-06 (Existencias y alertas):} para cada par \'unico
        \texttt{(id\_almacen, id\_producto)} se almacena \texttt{id\_sede},
        \texttt{existencia\_\allowbreak actual >= 0} y
        \texttt{stock\_minimo >= 0}
        (agregaci\'on \texttt{A2}). Cuando \texttt{existencia\_actual = 0} o
        \texttt{existencia\_actual < stock\_minimo}, el producto genera alerta
        de reabastecimiento.
  \item \textbf{RN-07 (Transferencias entre bodegas):} el almac\'en origen y
        el almac\'en destino deben ser distintos
        (\texttt{id\_almacen\_origen <> id\_almacen\_destino}). Las fechas
        siguen el orden l\'ogico \texttt{fecha\_despacho >= fecha\_solicitud} y
        \texttt{fecha\_recepcion >= fecha\_despacho} (admitiendo nulos en
        estados previos al despacho o recepci\'on). Cada l\'inea exige
        \texttt{cantidad\_enviada > 0} y permite \texttt{cantidad\_recibida}
        nula mientras el traslado est\'e en tr\'ansito, o \texttt{>= 0} al
        confirmarse en destino.
\end{itemize}

\subsection{Ventas, pagos y devoluciones (\texttt{ventas\_pagos})}
\begin{itemize}[leftmargin=1.5em, itemsep=2.5pt]
  \item \textbf{RN-08 (Pedidos de venta):} cada pedido vincula sede, almac\'en
        de despacho, cliente y vendedor (\texttt{id\_\allowbreak empleado}),
        con estado en \texttt{\{'Pendiente', 'Completado', 'Cancelado'\}} y
        \texttt{total\_registrado >= 0}. Cada producto aparece a lo sumo una vez
        por pedido en \texttt{detalle\_pedido} (agregaci\'on \texttt{A3}) con
        \texttt{cantidad > 0} y \texttt{precio\_unitario >= 0}.
  \item \textbf{RN-09 (Pagos y saldos):} todo pago registrado tiene
        \texttt{valor > 0}, medio permitido (\texttt{'Efectivo'},
        \texttt{'Tarjeta Debito'}, \texttt{'Tarjeta Credito'},
        \texttt{'Transferencia'}) y estado (\texttt{'Confirmado'},
        \texttt{'Rechazado'}, \texttt{'Pendiente'}). El saldo pendiente de un
        pedido equivale a \texttt{total\_registrado} menos la suma de sus pagos
        en estado \texttt{'Confirmado'}.
  \item \textbf{RN-10 (Devoluciones):} toda devoluci\'on se asocia a un pedido
        existente, con motivo obligatorio, estado en \texttt{\{'Solicitada',
        'Aprobada', 'Rechazada'\}} y l\'ineas con \texttt{cantidad > 0} y
        condici\'on (\texttt{'Buen Estado'}, \texttt{'Defectuoso'},
        \texttt{'Da\~nado'}).
\end{itemize}

\subsection{Identidad, seguridad y auditor\'ia (\texttt{identidad\_accesos} y \texttt{reportes})}
\begin{itemize}[leftmargin=1.5em, itemsep=2.5pt]
  \item \textbf{RN-11 (Cuentas 1:1 y roles):} cada empleado puede tener como
        m\'aximo un usuario del sistema (\texttt{id\_\allowbreak empleado} \'unico
        en \texttt{usuarios}) y un \texttt{nombre\_\allowbreak usuario} \'unico,
        con uno de los seis roles v\'alidos (\texttt{Administrador},
        \texttt{Gerencia}, \texttt{Inventario}, \texttt{Compras},
        \texttt{Vendedor}, \texttt{Auditor}).
  \item \textbf{RN-12 (Aislamiento y auditor\'ia inmutable):} cada rol de
        negocio solo modifica tablas del esquema de su responsabilidad. La tabla
        \texttt{auditoria\_accesos} no permite operaciones \texttt{UPDATE} ni
        \texttt{DELETE} por parte de roles operativos; \texttt{reportes} expone
        \'unicamente vistas \texttt{SELECT}.
\end{itemize}

% ============================================================================
\section{Arquitectura de microservicios y estrategia de distribuci\'on}

\subsection{Dominios y responsabilidad por esquema}

La base de datos se divide en seis esquemas independientes, detallados
gr\'aficamente en \texttt{Diagrama\_de\_\allowbreak microservicios\_BD.pdf}
y cuya matriz de permisos consta en \texttt{Matriz\_roles\_BD.pdf}:

\begin{center}
\renewcommand{\arraystretch}{1.32}
\begin{tabularx}{\linewidth}{@{}
    >{\ttfamily\footnotesize}p{4.5cm}
    >{\ttfamily\footnotesize}p{3.1cm}
    >{\centering\small\arraybackslash}p{1.4cm}
    >{\small\raggedright\arraybackslash}X @{}}
\toprule
\rowcolor{fondo}
{\sffamily\bfseries\footnotesize Esquema / Servicio} &
{\sffamily\bfseries\footnotesize Rol propietario} &
{\sffamily\bfseries\footnotesize Objetos} &
{\sffamily\bfseries\footnotesize Responsabilidad l\'ogica e interacci\'on} \\
\midrule
catalogo\_terceros & rol\_admin & 7 tablas &
  Maestros de sedes, almacenes, categor\'ias, productos, proveedores y clientes;
  consultado por compras, inventarios y ventas. \\
compras\_abastecimiento & rol\_compras & 4 tablas &
  Ciclo de \'ordenes de compra y recepciones; genera entradas de stock hacia
  inventarios. \\
inventarios\_transferencias & rol\_inventario & 3 tablas &
  Control de existencias por sede y almac\'en y traslados entre bodegas;
  atiende reservas y descuentos de ventas. \\
ventas\_pagos & rol\_vendedor & 5 tablas &
  Gesti\'on de pedidos de venta, recaudos y devoluciones por sede. \\
identidad\_accesos & rol\_admin & 3 tablas &
  Personal, credenciales de usuario y bit\'acora de accesos transversales. \\
reportes & {\sffamily\small Gerencia\,/\,Auditor} & 10 vistas &
  Capa de lectura consolidada (\texttt{SELECT}) sin tablas propias; integra
  los cinco esquemas operativos. \\
\bottomrule
\end{tabularx}
\end{center}

\subsection{Estrategia de fragmentaci\'on horizontal por sede}

Para responder tanto a consultas locales de cada sucursal como a consolidados
corporativos, el dise\~no define \textbf{fragmentaci\'on horizontal primaria
por sede} (criterio \texttt{LIST} sobre el atributo \texttt{id\_sede}) en dos
conjuntos operacionales de mayor volumen y localidad:

\begin{enumerate}[leftmargin=1.5em, itemsep=3pt]
  \item \textbf{\texttt{inventarios\_transferencias.inventarios}:} agrupa las
        existencias seg\'un la sede propietaria del almac\'en.
  \item \textbf{\texttt{ventas\_pagos.pedidos\_venta}:} agrupa los pedidos de
        venta seg\'un la sede donde se origina la operaci\'on.
\end{enumerate}

\noindent\textbf{Propiedades de la fragmentaci\'on:}
\begin{itemize}[leftmargin=1.5em, itemsep=2.5pt]
  \item \textbf{Condici\'on de fragmentaci\'on:} predicado de igualdad
        $F_k:\ \texttt{id\_sede} = k$ para $k \in \{1, 2, 3, 4, 5\}$ (cinco
        fragmentos correspondientes a las cinco sedes registradas en
        \texttt{catalogo\_terceros.sedes}).
  \item \textbf{Disjunci\'on y completitud:} como el atributo \texttt{id\_sede}
        es obligatorio (\texttt{NOT NULL}) y monovaluado en cada fila, cada
        registro pertenece exactamente a un \'unico fragmento (cero duplicaci\'on
        entre sedes) y la uni\'on de todos los fragmentos reconstruye el 100\,\%
        de la tabla l\'ogica.
  \item \textbf{Estado en Entrega 2 y materializaci\'on en Entrega 3:} en los
        scripts base actuales (\texttt{Scripts sql/\allowbreak
        creacion\_de\_tablas.sql}), ambas tablas se crean como tablas
        relacionales est\'andar que ya incluyen \texttt{id\_sede NOT NULL} y
        las claves \texttt{(id\_almacen, id\_producto)} y \texttt{(id\_pedido)}
        para verificar integridad y carga reproducible. En la fase de
        distribuci\'on (Entrega 3), los fragmentos se materializan mediante
        tablas f\'isicas por sede o particionamiento declarativo
        \texttt{PARTITION BY LIST (id\_sede)} (incluyendo \texttt{id\_sede} en
        la clave primaria compuesta como exige PostgreSQL para particiones
        nativas), integr\'andose en el esquema \texttt{reportes} a trav\'es de
        vistas con \texttt{UNION ALL} (\texttt{vista\_inventario\_\allowbreak
        por\_sede} y \texttt{vista\_ventas\_\allowbreak mensuales\_por\_sede}).
\end{itemize}

\subsection{Vistas de integraci\'on en \texttt{reportes}}

El esquema \texttt{reportes} contempla diez vistas de solo lectura para consumo
directo de gerencia, auditor\'ia y roles autorizados:
\texttt{vista\_\allowbreak inventario\_\allowbreak por\_sede},
\texttt{vista\_\allowbreak alertas\_\allowbreak reabastecimiento},
\texttt{vista\_\allowbreak productos\_\allowbreak sin\_inventario},
\texttt{vista\_\allowbreak ventas\_\allowbreak mensuales\_por\_sede},
\texttt{vista\_\allowbreak pedidos\_\allowbreak pendientes},
\texttt{vista\_\allowbreak saldos\_\allowbreak pendientes},
\texttt{vista\_\allowbreak compras\_\allowbreak pendientes},
\texttt{vista\_\allowbreak transferencias\_\allowbreak en\_transito},
\texttt{vista\_\allowbreak devoluciones\_\allowbreak por\_producto} y
\texttt{vista\_\allowbreak reporte\_\allowbreak corporativo} (con CTE agregadas por sede
para evitar productos cartesianos entre m\'ultiples relaciones 1:N).

% ============================================================================
\section{Modelo relacional y normalizaci\'on hasta 3FN}

\subsection{Esquema relacional completo (22 relaciones)}

En la notaci\'on siguiente, la clave primaria aparece \rpk{subrayada}, las
claves for\'aneas se se\~nalan con \textsuperscript{\scriptsize\sffamily FK}
y las restricciones de unicidad adicionales con
\textsuperscript{\scriptsize\sffamily UQ}. Los detalles completos de tipos de
datos y restricciones \texttt{CHECK} se encuentran en
\texttt{Diccionario\_de\_datos\_BD.pdf}.

\subsection*{Esquema \texttt{catalogo\_terceros}}
\rel{sedes}{\rpk{id\_sede}, \ruq{nombre\_sede}, ciudad, direccion, tipo\_sede, estado}
\rel{almacenes}{\rpk{id\_almacen}, \rfk{id\_sede}, nombre\_almacen, tipo\_almacen, estado}
\rel{categorias}{\rpk{id\_categoria}, \ruq{nombre\_categoria}, descripcion, estado}
\rel{productos}{\rpk{id\_producto}, \ruq{sku}, nombre, \rfk{id\_categoria}, unidad\_medida, precio\_lista, estado}
\rel{proveedores}{\rpk{id\_proveedor}, \ruq{identificacion\_proveedor}, razon\_social, ciudad\_proveedor, correo\_proveedor, telefono\_proveedor, estado}
\rel{producto\_proveedor}{\rpk{\rfk{id\_producto}}, \rpk{\rfk{id\_proveedor}}, costo\_referencia, plazo\_entrega}
\rel{clientes}{\rpk{id\_cliente}, tipo\_documento, \ruq{num\_documento}, nombre\_completo\_cliente, correo\_cliente, telefono\_cliente, estado}

\subsection*{Esquema \texttt{compras\_abastecimiento}}
\rel{ordenes\_compra}{\rpk{id\_orden\_compra}, \rfk{id\_proveedor}, \rfk{id\_almacen\_destino}, \rfk{id\_empleado}, fecha\_orden, fecha\_estimada, estado}
\rel{detalle\_orden\_compra}{\rpk{\rfk{id\_orden\_compra}}, \rpk{\rfk{id\_producto}}, cantidad\_solicitada, costo\_unitario}
\rel{recepciones\_compra}{\rpk{id\_recepcion}, \rfk{id\_orden\_compra}, fecha\_recepcion, estado}
\rel{detalle\_recepcion}{\rpk{\rfk{id\_recepcion}}, \rpk{\rfk{id\_producto}}, cantidad\_recibida}

\subsection*{Esquema \texttt{inventarios\_transferencias}}
\rel{inventarios}{\rfk{id\_sede}, \rpk{\rfk{id\_almacen}}, \rpk{\rfk{id\_producto}}, existencia\_actual, stock\_minimo, fecha\_actualizacion}
\rel{transferencias}{\rpk{id\_transferencia}, \rfk{id\_almacen\_origen}, \rfk{id\_almacen\_destino}, fecha\_solicitud, fecha\_despacho, fecha\_recepcion, estado}
\rel{detalle\_transferencia}{\rpk{\rfk{id\_transferencia}}, \rpk{\rfk{id\_producto}}, cantidad\_enviada, cantidad\_recibida}

\subsection*{Esquema \texttt{ventas\_pagos}}
\rel{pedidos\_venta}{\rpk{id\_pedido}, \rfk{id\_sede}, \rfk{id\_almacen}, \rfk{id\_cliente}, \rfk{id\_empleado}, fecha\_pedido, estado, total\_registrado}
\rel{detalle\_pedido}{\rpk{\rfk{id\_pedido}}, \rpk{\rfk{id\_producto}}, cantidad, precio\_unitario}
\rel{pagos}{\rpk{id\_pago}, \rfk{id\_pedido}, fecha\_pago, valor, medio\_pago, estado}
\rel{devoluciones}{\rpk{id\_devolucion}, \rfk{id\_pedido}, fecha\_solicitud, motivo, estado}
\rel{detalle\_devolucion}{\rpk{\rfk{id\_devolucion}}, \rpk{\rfk{id\_producto}}, cantidad, condicion\_producto}

\subsection*{Esquema \texttt{identidad\_accesos}}
\rel{empleados}{\rpk{id\_empleado}, \rfk{id\_sede}, \ruq{num\_documento}, nombre\_completo, cargo, correo, estado}
\rel{usuarios}{\rpk{id\_usuario}, \rfk{\ruq{id\_empleado}}, \ruq{nombre\_usuario}, password\_hash, rol, estado, fecha\_creacion}
\rel{auditoria\_accesos}{\rpk{id\_auditoria}, \rfk{id\_usuario}, fecha\_hora, accion, objeto\_consultado}

\subsection{Justificaci\'on de normalizaci\'on hasta Tercera Forma Normal (3FN)}

\begin{enumerate}[leftmargin=1.5em, itemsep=4pt]
  \item \textbf{Primera Forma Normal (1FN):} todas las tablas poseen clave
        primaria definida, dominios at\'omicos (por ejemplo, \texttt{tipo\_documento}
        separado de \texttt{num\_documento}) y ausencia de grupos repetitivos.
        Los \'items m\'ultiples de \'ordenes, recepciones, transferencias,
        pedidos y devoluciones se extraen a sus respectivas tablas de detalle.

  \item \textbf{Segunda Forma Normal (2FN):} las 15 entidades con clave primaria
        simple cumplen 2FN autom\'aticamente al no existir dependencias
        parciales. En las \textbf{7 tablas de clave primaria compuesta}
        (\texttt{producto\_\allowbreak proveedor},
        \texttt{detalle\_\allowbreak orden\_\allowbreak compra},
        \texttt{detalle\_\allowbreak recepcion}, \texttt{inventarios},
        \texttt{detalle\_\allowbreak transferencia},
        \texttt{detalle\_\allowbreak pedido} y
        \texttt{detalle\_\allowbreak devolucion}), cada atributo no clave depende
        funcionalmente de la clave completa: por ejemplo,
        \texttt{costo\_referencia} y \texttt{plazo\_entrega} dependen de la
        combinaci\'on \texttt{(id\_producto, id\_proveedor)};
        \texttt{cantidad} y \texttt{precio\_unitario} dependen del par
        \texttt{(id\_pedido, id\_producto)}; y \texttt{existencia\_\allowbreak actual} y
        \texttt{stock\_minimo} dependen del par
        \texttt{(id\_almacen, id\_producto)}.

  \item \textbf{Tercera Forma Normal (3FN) y desnormalizaci\'on controlada:}
        los atributos descriptivos de sedes, categor\'ias, proveedores,
        clientes y empleados residen \'unicamente en sus entidades maestras,
        evitando dependencias transitivas en las tablas transaccionales. Se
        mantienen tres decisiones conscientes de dise\~no que no corresponden a
        anomal\'ias sino a requisitos operativos y de distribuci\'on:
        \begin{itemize}[leftmargin=1.3em, itemsep=2pt]
          \item \texttt{inventarios.id\_sede} y \texttt{pedidos\_venta.id\_sede}:
                aunque la sede puede inferirse a trav\'es del almac\'en, ambos
                atributos se incluyen expl\'icitamente como \textbf{clave de
                fragmentaci\'on horizontal} para enrutar particiones por sede
                sin consultar el esquema \texttt{catalogo\_terceros}.
          \item \texttt{pedidos\_\allowbreak venta.\allowbreak total\_\allowbreak registrado}:
                campo calculado como la suma de \texttt{cantidad} $\times$
                \texttt{precio\_\allowbreak unitario} en
                \texttt{detalle\_pedido}, mantenido actualizado para agilizar
                la consulta de saldos y reportes corporativos.
          \item \texttt{detalle\_pedido.\allowbreak precio\_unitario} y
                \texttt{detalle\_\allowbreak orden\_\allowbreak compra.\allowbreak
                costo\_\allowbreak unitario}:
                representan el valor hist\'orico pactado en la fecha de la
                transacci\'on y no una copia redundante del precio actual del
                cat\'alogo.
        \end{itemize}
\end{enumerate}

\end{document}